

Supposons @@M0@@. Alors @@M1@@ ; donc @@M2@@ et @@M3@@. De plus, @@M4@@ ; donc @@M5@@. Ainsi, @@M6@@ est bien un @@M7@@-faisceau. En outre, 6.5 (ii) montre que @@M8@@. Si @@M9@@, on peut supposer @@M10@@ un @@M11@@-faisceau et @@M12@@, en raisonnant par r\'ecurrence sur @@M13@@. Alors 6.5 montre que @@M14@@ et @@M15@@ ; d'o\`u l'assertion.

\textbf{Th\'eor\`eme 6.7.} \emph{Il existe deux suites @@M16@@ et @@M17@@ de polyn\^omes ayant les propri\'et\'es suivantes. On suppose les conditions de 6.1 remplies et @@M18@@ un quotient de @@M19@@ avec @@M20@@.}
\begin{enumerate}[label=(\roman*)]
\item \emph{Si @@M21@@ est un @@M22@@-faisceau, alors pour @@M23@@,}
@@M24@@
\item \emph{Si @@M25@@ est le polyn\^ome de Hilbert de @@M26@@, alors pour @@M27@@,}
@@M28@@
\end{enumerate}
\emph{(Remarquons que @@M29@@ et @@M30@@ puisque @@M31@@).}

En effet, soit une suite exacte @@M32@@ telle que @@M33@@ soit un @@M34@@-faisceau ; raisonnant par r\'ecurrence sur @@M35@@, on peut supposer d\'efinis @@M36@@, compte tenu de 1.7. Alors 6.4 (ii) met en \'evidence la d\'efinition de @@M37@@, compte tenu de 1.8. Enfin (ii) r\'esulte formellement de 6.4 (iii), 6.6, et l'assertion (i) (o\`u on prend @@M38@@ pour @@M39@@).

